%
%	Configure LaTeX to produce a PDF presentation using the beamer class
%


% Default is 16:9 widescreen ratio
\documentclass[ignorenonframetext,12pt,aspectratio=169]{beamer}


% Create a handout with the slides on a page
% \documentclass[handout]{beamer}
% \usepackage{pgfpages}
% \usepackage{handoutWithNotes}
% % \pgfpagesuselayout{4 on 1}[letterpaper,landscape,border shrink=5mm]
% \pgfpagesuselayout{3 on 1 with notes}




% Allow things like '<a>' in verbatim environments (fixes font encoding issue)
\usepackage[T1]{fontenc}


% Minimum required packages to suppport MMD features
\usepackage[sort&compress]{natbib}


% For each part, show a title page and an outline of that part
\AtBeginPart
{
   \begin{frame}
      \partpage
   \end{frame}

   \begin{frame}
       \frametitle{Outline}
       \tableofcontents
   \end{frame}
}

% For each section, show an outline highlighting where we are
\AtBeginSection[]
{
   \begin{frame}
       \frametitle{Outline}
       \tableofcontents[currentsection,currentsubsection]
   \end{frame}
}
\def\mytitle{Sample Beamer Presentation}
\def\mysubtitle{With a subtitle, no less!}
\def\myauthor{Fletcher T. Penney}
\def\affiliation{http://fletcherpenney.net/multimarkdown}
\def\mycopyright{2026 Fletcher T. Penney This document is licensed under a Creative Commons License. http://creativecommons.org/licenses/by-sa/2.5/}
%
%	Get ready for the actual document
%


% Use default MMD metadata for beamer equivalents

\ifx\mysubtitle\undefined
\else
	\subtitle{\mysubtitle}
\fi

\ifx\affiliation\undefined
\else
	\institute{\affiliation}
\fi

\ifx\mydate\undefined
	\def\mydate{\today}
\else
	\date{\mydate}
\fi


% This does not look good in all themes
\ifx\event\undefined
\else
	\date[\mydate]{\mydate~ / \event }
\fi


% Show "current/total" slide counter in footer
% Works differently in different themes
\title[\mytitle\hspace{2em}\insertframenumber/
\inserttotalframenumber]{\mytitle}

% Set author
\author{\myauthor}

% Include space between paragraphs
\addtolength{\parskip}{\baselineskip}

\ifx\theme\undefined
\else
	\usetheme{\theme}
\fi

\begin{document}

\maketitle

\part{Part 1}
\label{part1}

\section{Section A}
\label{sectiona}

\subsection{Subsection 1}
\label{subsection1}

\begin{frame}[fragile]
\frametitle{MultiMarkdown}
\label{multimarkdown}

Foo.

\end{frame}

\begin{frame}[fragile]
\frametitle{More MultiMarkdown}
\label{moremultimarkdown}

Foo.

\end{frame}

\subsection{Subsection 2}
\label{subsection2}

\begin{frame}[fragile]
\frametitle{Still MultiMarkdown}
\label{stillmultimarkdown}

Foo.

\end{frame}

\section{Section B}
\label{sectionb}

\subsection{Subsection a}
\label{subsectiona}

\begin{frame}[fragile]
\frametitle{Again}
\label{again}

Bar.

\end{frame}

\subsection{Subsection b}
\label{subsectionb}

\begin{frame}[fragile]
\frametitle{Yet Again}
\label{yetagain}

Bar.

\end{frame}

\part{Part 2}
\label{part2}

\section{Section 2A}
\label{section2a}

\subsection{SubSection 2A}
\label{subsection2a}

\begin{frame}[fragile]
\frametitle{Wrap it up}
\label{wrapitup}

Baz.
\end{frame}

\mode<all>
% Durham Theme has a thank you slide
\ifdefined\makethankyou
\makethankyou
%\makethankyou{Foo}{Bar}
\fi
\end{document}\mode*
